\cotefile{157-3}
\scanpage{66}
\mspage{15}

et soit \[(f_i)_{i\in I}:(X_i)_{i\in I}\longrightarrow(Y_i)_{i\in I}\] un \unsure{hom} de systèmes inductifs dans $\mathrm{Cat}$, indexés par un ens. ordonné filtrant, soient $X=\varinjlim X_i$, $Y=\varinjlim Y_i$, $f=\varinjlim f_i$, \deleted{$f:\varinjlim_i X_i\,\unsure{\to|=}$} $f:X\longrightarrow Y$. Si $f_i$ est dans $W$ pour tt $i$, alors $f\in W$.

On factorise cha\uline{que} $f_i$ en \[X_i\xrightarrow{\alpha_i}\overline X_i\xrightarrow{\overline f_i}Y_i,\qquad \alpha_i\in W_{\unsure{\omega|\infty}},\quad\overline f_i\in\mathrm{Prof}_{W_\omega}\]
% slanted
\margin{$W_\omega$ introduit dans la condition $4\,\mathrm b''$, page 14} \unsure{suivant} le système \unsure{décrit}. La construction étant fonctorielle, les $\overline X_i$ forment eux aussi un syst. inductif. De plus, comme le \deleted{\illegible{1}} foncteur \deleted{\illegible{1}} $S\longrightarrow\underline{\mathrm{Ch}}(S)$ qui donne naissance \deleted{\unsure{au}} à la factorisation canonique des morphismes, commute aux $\varinjlim$ filtrantes (immédiat), il s'ensuit que, si on pose \[\overline X=\varinjlim_i\overline X_{\unsure{i}},\] et qu'on regarde les flèches \[X\xrightarrow{\alpha}\overline X\xrightarrow{\overline f}Y,\qquad \alpha=\varinjlim\alpha_i,\quad\overline f=\varinjlim\overline f_i,\] avec $\overline f\alpha=f$, cette factorisation de $f$ est bien sa factorisation canonique. \inserted{Donc $\alpha\in W$, $\overline f\in\mathrm{Prof}_{W_\omega}$.} Ceci vu, en vertu des propriétés \deleted{\illegible{1}} $f\in W$, \deleted{\unsure{équivaut à dire que}} \inserted{i.e. $\overline f\in W$,} est vrai si les fibres de $\overline f$ \deleted{\unsure{sont}} \unsure{soient|sont} $W$-asphériques. Mais les fibres de $\overline f$ s'interprètent comme limites inductives filtrantes de fibres des $\overline f_i$. Lesquelles sont $W$-asphériques par \deleted{$4\,\unsure{\mathrm b}$} la condition $4\,\mathrm b''$ des fibrations. Donc l'axiome \textcircled{6} donne ce qu'il faut.

% slanted
\margin{et peut-être \textcircled{2} sous la forme plus \unsure{faible} \textcircled{\scriptsize 2bis}?} \uline{Conjecture} (\unsure{hardie}? \unsure{folle}?) Si $W\subset\mathrm{Fl}(\mathrm{Cat})$ satisfait les axiomes \textcircled{1} à \textcircled{6} (avec 4 sous sa forme la plus faible $4\,\mathrm b''$), alors $W=W_\infty$.

% archive sheet number: 66
